\documentclass[11pt]{article}
\usepackage[a4paper,margin=1in]{geometry}
\usepackage{enumitem}
\usepackage{parskip}
\usepackage{microtype}
\usepackage{titlesec}

\titlespacing*{\section}{0pt}{1.4ex plus .2ex}{0.6ex}
\titleformat{\section}{\normalfont\large\bfseries}{}{0em}{}
\setlength{\parskip}{0.5em}
\setlist[enumerate]{leftmargin=1.4em,topsep=3pt,itemsep=4pt}

\title{\vspace{-3em}Data Proposal\\[0.2em]
\large Item-response data for measuring ability and tracking its change}
\author{}
\date{}

\begin{document}
\maketitle
\vspace{-3.5em}

\section*{Research context}

We are developing a measurement model that estimates a respondent's
ability from a history of item responses and tracks how that ability
changes over time. It places different response formats on a single
common scale and can link new items onto an existing scale without a full
recalibration. We have evaluated it on simulated and public data, and we
are now seeking real response data to test it at scale. We would value
whatever can be shared, suitably anonymized and under an agreement that
works for you.

\section*{Data of interest}

Listed roughly in order of value to us, and any single category is useful
on its own. In every case a respondent identifier, an item or test
identifier, a score, and a date or sequence index are usually enough to
begin. Richer detail is welcome but not required.

\begin{enumerate}
  \item \textbf{Longitudinal response data.} Repeated measurements of the
  same respondents across occasions, terms, or years. This is the most
  valuable category, since it lets us estimate ability growth and its
  rate.

  \item \textbf{Item-level responses.} Scored responses at the item level
  rather than test totals, with a reasonable number of respondents per
  item. Dichotomous and polytomous scoring are equally welcome.

  \item \textbf{Existing item analyses.} Any item parameters or
  calibrations already available for these items, together with the model
  used. These give us a reference to benchmark our own estimates against.

  \item \textbf{Linking or equating designs.} Response data connected
  across forms, groups, or years through common items. Knowing which items
  are shared is the only addition we need here. This supports our work on
  extending a scale without recalibration.

  \item \textbf{Mixed-format assessments.} Tests combining formats, for
  example selected-response with constructed-response or rubric-scored
  items, ideally for the same respondents. A format label per item is
  helpful. These let us place several formats on one scale.

  \item \textbf{Item content or descriptors (optional).} Item text, or
  topic and skill classifications, carrying no respondent information.
  This supports content-based placement of new items and the modeling of
  skill dimensions.
\end{enumerate}

\section*{Scope, format, and governance}

A small extract is enough to begin. One subject area, a few thousand
respondents, or a limited time window would already let us run a first
study. We can work with anonymized or pseudonymized data in any common
tabular format, at whatever level of detail and coverage you are
comfortable releasing. We are glad to work under a data-sharing or
non-disclosure agreement, to restrict use to research, and not to
redistribute. A small pilot first, then more if it proves useful, suits us
well.

\end{document}
